% Self-contained retraction reduction from scalar UR to zero-equation induction.
\section{Zero-equation induction from uniqueness}
\label{sec:ur-retraction-complete}

The argument in this section is the retraction argument already present in
the incoming \texttt{scope\_kernel} report.  It is reproduced to make the
dependency on equality reflection exact.  It uses neither predecessor nor
a general numerical equality characteristic.

All calculations are under one fixed retained first-order context
$\Gamma$.  Fresh displayed variables are absent from $\Gamma$ and all
passive parameters.  UR first gives
\[
                         I(n,0;u.Su)=n.                  \tag{Id}
\]
Use the template $f(z)=z$ and update $h(u)=Su$; its recurrence premise is
reflexivity.  Common-update uniqueness, obtained by applying UR separately
to two terms and replacing equal seeds, is used below.

\subsection{Retraction rigidity}
Let $e(x)$ be a fixed scalar term template.  Suppose $\Gamma$ proves
\[
 e(0)=0,\qquad e(e(x))=e(x),\qquad e(S(e(x)))=S(e(x)).    \tag{Ret}
\]
Then $\Gamma$ proves $e(t)=t$ for every target $t$.

To see this, put $h(u)=e(S(e(u)))$ and $r(n)=I(n,0;u.h(u))$.
Idempotence in Ret gives $e(h(u))=h(u)$ and $h(e(u))=h(u)$.
Therefore the template $f(n)=e(r(n))$ has initial value zero and
unconditional recurrence
\[
 f(Sn)=e(h(r(n)))=h(r(n))=h(e(r(n)))=h(f(n)).
\]
UR gives $e(r(n))=I(n,0;u.h(u))=r(n)$.
The last equation of Ret now yields
\[
 r(Sn)=h(r(n))=S(e(r(n)))=S(r(n)),\qquad r(0)=0.
\]
A second application of UR, followed by Id, gives $r(n)=n$.
Replacing $r(n)$ by $n$ in the previous equality proves $e(n)=n$.
No induction was used, and no claim that the standard finite successor
orbit exhausts a model was used.

\subsection{A numerical split suffices for induction}
For a fixed formula $A(n)$, suppose
\[
 \Gamma\vdash A(0),\qquad \Gamma,A(z)\vdash A(Sz).
\]
Suppose moreover that $A(n)$ is decidable under $\Gamma$ and that a term
$e(n)$ satisfies $e(n)=n$ under $A(n)$ and $e(n)=0$ under $\neg A(n)$.
Cases give $A(e(n))$: in the first case use $A(n)$, and in the second
use the proved base $A(0)$.  Applying the true-case identity at $e(n)$
gives $e(e(n))=e(n)$.  The induction-step derivation, instantiated at
$e(n)$, gives $A(S(e(n)))$, hence $e(S(e(n)))=S(e(n))$.
The base similarly gives $e(0)=0$.  Retraction rigidity proves $e(n)=n$;
transporting $A(e(n))$ gives $A(n)$.

This is a metatheorem about a justified numerical split.  It supplies no
decision procedure, comprehension principle, or induction schema for
arbitrary quantified formulas.

\subsection{The unconditional zero-equation case}
Write $1=S0$ and $C(v,a,b)=I(v,a;u.b)$.  The two computations give
$C(0,a,b)=a$ and $C(Sv,a,b)=b$.
The separation axiom $0=1\to\bot$ implies $Sv\ne0$: otherwise applying
$w\mapsto C(w,0,1)$ gives $1=0$, hence a contradiction.
Coverage therefore decides $q(n)=0$ for every fixed scalar term $q(n)$.
The term
\[
                             e(n)=C(q(n),n,0)
\]
has exactly the split required in the preceding subsection.
Consequently the following contextual rule is derived:
\[
 \frac{\Gamma\vdash q(0)=0\qquad
       \Gamma,q(z)=0\vdash q(Sz)=0}
      {\Gamma\vdash q(t)=0}.                            \tag{ZI}
\]
Every substitution and case distinction is intuitionistic.  The retained
context can contain arbitrary quantified formulas, provided the fresh
induction variable is absent from it.
